\section{A weighted prime sum and its Laplace transform}\label{sec:transform}

The argument does not bound a sum over primes with sign changes. It bounds from above a sum with
nonnegative weights, and the zeros of the Dirichlet $L$-functions are seen in the poles of its
Laplace transform. This section defines the sum, computes the transform and its residues, and
records the elementary bounds that the later sections use.

\subsection{The weights}\label{ss:weights}
Throughout this section $d$ is a prime and $\chiK=(\,\cdot\,/d)$ is the Legendre symbol, extended
to $\Z$ by $\chiK(n)=0$ for $d\mid n$. For $d\equiv 3\pmod 4$ it is the quadratic character of
$K=\Q(\sqrt{-d})$; the letter $K$ anticipates this case. Let $q\ge 1$ be an integer. For a
Dirichlet character $\psi$ modulo $q$ we write $\psi\chiK$ for the character modulo $qd$ given by
$n\mapsto\psi(n)\chiK(n)$, and we put
\[
  E_\psi(s)=
  \begin{cases}
    L(s,\psi) & \text{if $\psi$ is not principal,}\\
    (s-1)L(s,\psi) & \text{if $\psi$ is principal.}
  \end{cases}
\]
In both cases $E_\psi$ is entire, and $E_\psi(s)\neq0$ for $\re s\ge1$.

\begin{definition}\label{def:lambda}
For $n\ge1$ let
\[
  \lambda_q(n)=
  \begin{cases}
    \Lam(n)\bigl(1-\chiK(n)\bigr) & \text{if } n\equiv 1\pmod q,\\
    0 & \text{otherwise,}
  \end{cases}
\]
where $\Lam$ is the von Mangoldt function. Let
\[
  G_q(s)=\frac{1}{\totient(q)}\Bigl(\sum_{\psi\bmod q}
    \Bigl(\frac{E_{\psi\chiK}'}{E_{\psi\chiK}}(s)-\frac{E_\psi'}{E_\psi}(s)\Bigr)-1\Bigr),
\]
the sum running over all Dirichlet characters modulo $q$, and let
\begin{equation}\label{eq:Rq}
  R_q(x)=\sum_{1\le n\le x}\lambda_q(n)\Bigl(\frac1n-\frac1x\Bigr)-\frac{\log x}{\totient(q)}-c_q,
  \qquad c_q=\re G_q(1),
\end{equation}
for real $x>0$. We call $R_q$ the \emph{remainder} modulo $q$.
\end{definition}
\lean{New.lambdaKResidueClass, New.logDerivMean, New.primeRemainder}

The weight $\lambda_q(n)$ vanishes at the prime powers $n=p^k$ with $\chiK(p^k)=1$; in particular,
for $d\equiv 3\pmod 4$, at all powers of primes that split in $K$ and at the even powers of inert
primes. The function $G_q$ is meromorphic, with poles at most at the zeros of the $E_\psi$ and the
$E_{\psi\chiK}$. By \cref{lem:G-real}, $G_q(1)$ is real when $d$ is odd and prime
to $q$, so that $c_q=G_q(1)$; and in \cref{thm:pnt-Rq} that $c_q$ is then the constant for which
$R_q(x)\to0$.

\begin{lemma}\label{lem:lambda-bounds}
For all $n\ge1$ one has $0\le\lambda_q(n)\le2\Lam(n)\le2\log n\le 2n$, and for every integer
$N\ge1$
\begin{equation}\label{eq:mertens-two}
  \sum_{1\le n\le N}\frac{\Lam(n)}{n}\le 2\log N.
\end{equation}
\end{lemma}
\lean{New.lambdaKResidueClass_nonneg, New.lambdaKResidueClass_le_two_mul_vonMangoldt, New.lambdaKResidueClass_le_two_mul, New.sum_vonMangoldt_div_le}

\begin{proof}
Since $\chiK(n)\in\{-1,0,1\}$, the factor $1-\chiK(n)$ lies in $[0,2]$, which gives the first two
inequalities; the others are $\Lam(n)\le\log n\le n$. For \eqref{eq:mertens-two} we use
$N/n\le 2\lfloor N/n\rfloor$ for $1\le n\le N$ and the identity
$\sum_{n\le N}\Lam(n)\lfloor N/n\rfloor=\log N!$, which follows from $\sum_{e\mid k}\Lam(e)=\log k$.
Hence
\[
  N\sum_{n\le N}\frac{\Lam(n)}{n}\le 2\sum_{n\le N}\Lam(n)\Bigl\lfloor\frac Nn\Bigr\rfloor
  =2\log N!\le 2N\log N.\qedhere
\]
\end{proof}

For $r\in\{0,1\}$ and $y>0$ put
\[
  A_r(y)=y^{-r}\sum_{1\le n\le y}\lambda_q(n)\,n^{r-1}.
\]
Then $R_q(x)=A_0(x)-A_1(x)-\log x/\totient(q)-c_q$. The sum defining $A_r(y)$ has at most $y$ terms,
each at most $2y^r$ by \cref{lem:lambda-bounds}, so
\begin{equation}\label{eq:Ar-bound}
  \abs{A_r(e^v)}\le 2e^v\qquad(v\in\R,\ r\in\{0,1\}).
\end{equation}
Since $v/\totient(q)\le v\le e^v$ for $v\ge0$, it follows that
\begin{equation}\label{eq:R-growth}
  \abs{R_q(e^v)}\le(\abs{c_q}+5)\,e^v\qquad(v\ge0).
\end{equation}
\lean{New.abs_lambdaKSum_exp_le, New.abs_primeRemainder_exp_le}

We shall often use the following form of the definition. For $0<z\le A$ with $A$ an integer,
\begin{equation}\label{eq:R-posPart}
  R_q(z)=\frac1A\sum_{1\le n\le A}\lambda_q(n)\Bigl(\frac An-\frac Az\Bigr)_{\!+}
  -\frac{\log z}{\totient(q)}-c_q,
\end{equation}
where $t_+=\max(t,0)$: the terms with $n>z$ vanish, and for $n\le z$ the positive part is the
number itself.
\lean{New.primeRemainder_eq_sum_posPart, New.primeRemainder_eq_sum}

\subsection{The Dirichlet series and the transform}\label{ss:dirichlet-series}

\begin{lemma}\label{lem:twist-not-principal}
Let $d$ be an odd prime not dividing $q$. For every character $\psi$ modulo $q$ the character
$\psi\chiK$ modulo $qd$ is not principal.
\end{lemma}
\lean{New.quadTwist_ne_one}

\begin{proof}
Let $a$ be a quadratic non-residue modulo $d$, and choose $n$ with $n\equiv1\pmod q$ and
$n\equiv a\pmod d$. Then $(\psi\chiK)(n)=\psi(1)\chiK(a)=-1$, while a principal character takes
only the values $0$ and $1$.
\end{proof}

\begin{lemma}\label{lem:dirichlet-series}
Let $d$ be an odd prime not dividing $q$. For $\re s>1$,
\[
  \sum_{n\ge1}\lambda_q(n)\,n^{-s}=\frac{s}{\totient(q)(s-1)}+G_q(s).
\]
\end{lemma}
\lean{New.LSeries_lambdaKResidueClass_eq_add_logDerivMean}

\begin{proof}
An integer $n\equiv1\pmod q$ is prime to $q$, and for $n$ prime to $q$ orthogonality gives
$\sum_{\psi\bmod q}\psi(n)=\totient(q)$ or $0$ according as $n\equiv1\pmod q$ or not. Hence
\[
  \lambda_q(n)=\frac{1}{\totient(q)}\sum_{\psi\bmod q}
  \bigl(\psi(n)\Lam(n)-(\psi\chiK)(n)\Lam(n)\bigr)\qquad(n\ge1).
\]
For a character $\chi$ and $\re s>1$ we have $\sum_n\chi(n)\Lam(n)n^{-s}=-L'(s,\chi)/L(s,\chi)$,
which is $-E_\chi'/E_\chi(s)$ if $\chi$ is not principal and $-E_\chi'/E_\chi(s)+1/(s-1)$ if it is.
Among the characters $\psi$ and $\psi\chiK$ only the principal character modulo $q$ is principal,
by \cref{lem:twist-not-principal}. Summing,
\[
  \sum_{n\ge1}\lambda_q(n)n^{-s}
  =\frac{1}{\totient(q)}\Bigl(\frac{1}{s-1}
   +\sum_{\psi}\Bigl(\frac{E_{\psi\chiK}'}{E_{\psi\chiK}}(s)-\frac{E_\psi'}{E_\psi}(s)\Bigr)\Bigr)
  =\frac{1}{\totient(q)}\Bigl(\frac{1}{s-1}+1\Bigr)+G_q(s),
\]
and $1/(s-1)+1=s/(s-1)$.
\end{proof}

\begin{lemma}\label{lem:transform}
Let $d$ be an odd prime not dividing $q$. For $\re s>2$,
\[
  \int_0^\infty R_q(e^v)\,e^{(1-s)v}\dd v=\frac{G_q(s)/s-c_q}{s-1}.
\]
\end{lemma}
\lean{New.LSeries_lambdaKResidueClass, New.integral_primeRemainder}

\begin{proof}
Write $D(s)$ for the Dirichlet series of \cref{lem:dirichlet-series}. Partial summation applied
to the coefficients $\lambda_q(n)n^{r-1}$ at the point $s+r-1$, followed by the substitution
$t=e^v$, gives
\[
  D(s)=(s+r-1)\int_0^\infty A_r(e^v)\,e^{(1-s)v}\dd v\qquad(r\in\{0,1\},\ \re s>2);
\]
the integrals converge absolutely by \eqref{eq:Ar-bound}, and the partial sums of the coefficients
are $O(y^{1+r})$, which is what partial summation needs on $\re(s+r-1)>1+r$. Since
$\int_0^\infty v^k e^{(1-s)v}\dd v=k!/(s-1)^{k+1}$, we obtain
\[
  \int_0^\infty R_q(e^v)e^{(1-s)v}\dd v
  =\frac{D(s)}{s-1}-\frac{D(s)}{s}-\frac{1}{\totient(q)(s-1)^2}-\frac{c_q}{s-1}.
\]
By \cref{lem:dirichlet-series}, $D(s)\bigl(\frac1{s-1}-\frac1s\bigr)=\frac{D(s)}{s(s-1)}
=\frac{1}{\totient(q)(s-1)^2}+\frac{G_q(s)}{s(s-1)}$, and the assertion follows.
\end{proof}

\begin{lemma}\label{lem:G-real}
Let $d$ be an odd prime not dividing $q$. Then $G_q$ is holomorphic on $\re s\ge1$, and $G_q(1)$
is real; thus $c_q=G_q(1)$.
\end{lemma}
\lean{New.analyticAt_logDerivMean_of_one_le_re, New.ofReal_re_logDerivMean_one}

\begin{proof}
The functions $E_\psi$ and $E_{\psi\chiK}$ have no zeros on $\re s\ge1$, so $G_q$ is holomorphic
there, and $F(s)=G_q(s)/s-c_q$ is holomorphic on a neighbourhood of $[1,\infty)$. For real $x>2$,
\cref{lem:transform} expresses $F(x)$ as $(x-1)$ times the integral of a real function, so $F$ is
real on $(2,\infty)$. The function $x\mapsto\im F(x)$ is real analytic on $[1,\infty)$ and vanishes
on $(2,\infty)$, hence on $[1,\infty)$. At $x=1$ this says $\im G_q(1)=0$.
\end{proof}

\subsection{Residues}\label{ss:residues}
For a function $f$ meromorphic at $\rho$ we write $\ord_\rho f$ for the order of vanishing of $f$
at $\rho$ when $f$ is holomorphic there; all orders below are of entire functions, hence
nonnegative.

\begin{lemma}\label{lem:residue-G}
For every $\rho\in\C$,
\[
  \lim_{z\to\rho}(z-\rho)\,G_q(z)=\frac{1}{\totient(q)}\sum_{\psi\bmod q}
  \bigl(\ord_\rho E_{\psi\chiK}-\ord_\rho E_\psi\bigr).
\]
Consequently, for every constant $\kappa$ and every $\rho\neq0,1$,
\[
  \lim_{z\to\rho}(z-\rho)\,\frac{G_q(z)/z-\kappa}{z-1}
  =\frac{1}{\totient(q)\,\rho(1-\rho)}\sum_{\psi\bmod q}
  \bigl(\ord_\rho E_{\psi}-\ord_\rho E_{\psi\chiK}\bigr).
\]
\end{lemma}
\lean{New.tendsto_sub_mul_logDerivMean, New.tendsto_logDerivMean_residue}

\begin{proof}
For an entire function $f\not\equiv0$, $(z-\rho)f'(z)/f(z)\to\ord_\rho f$ as $z\to\rho$. Summing over
$\psi$ gives the first limit; the term $-1/\totient(q)$ contributes nothing since $z-\rho\to0$. The
second limit is the first divided by $\rho(\rho-1)$, the constant $\kappa$ again contributing
nothing.
\end{proof}

The transform of $R_q$ thus continues to a meromorphic function whose poles in the right half-plane
sit at zeros of the $E_\psi$ and the $E_{\psi\chiK}$. We record exactly what will be used.

\begin{definition}\label{def:polar}
Let $E\colon\R\to\R$ and $r\colon\C\to\C$. We say that $E$ is \emph{polar with residues $r$} if
$E$ is measurable, there are a constant $C$ with $\abs{E(v)}\le Ce^v$ for $v\ge0$ and a meromorphic
function $\Phi$ on $\C$ with
\[
  \Phi(s)=\int_0^\infty E(v)\,e^{(1-s)v}\dd v\qquad(\re s>2),
\]
and $(z-\rho)\Phi(z)\to r(\rho)$ as $z\to\rho$, $z\neq\rho$, for every $\rho$ with $\re\rho>0$.
\end{definition}
\lean{New.IsPolar}

In particular $\Phi$ has at most simple poles in $\re s>0$, and $r(\rho)=0$ where $\Phi$ is
holomorphic. For $E(v)=R(e^v)$ the integral is $\int_1^\infty R(y)y^{-s}\dd y$.

\begin{proposition}\label{prop:polar}
Let $d\equiv7\pmod8$ be a prime not dividing $q$. Then $v\mapsto R_q(e^v)$ is polar with residues
\[
  r_q(s)=\frac{1}{\totient(q)\,s(1-s)}\sum_{\psi\bmod q}
  \bigl(\ord_s E_\psi-\ord_s E_{\psi\chiK}\bigr)\quad(s\neq0,1),\qquad r_q(1)=0.
\]
\end{proposition}
\lean{New.isPolar_primeRemainder}

\begin{proof}
Measurability is clear, and the growth bound is \eqref{eq:R-growth}. Put
$\Phi(s)=(G_q(s)/s-c_q)/(s-1)$, a meromorphic function on $\C$; by \cref{lem:transform} it is the
transform for $\re s>2$. At $\rho\ne1$ with $\re\rho>0$ the residue is given by
\cref{lem:residue-G}. At $\rho=1$ the numerator $F(s)=G_q(s)/s-c_q$ is holomorphic and vanishes, by
\cref{lem:G-real}; so $(z-1)\Phi(z)=F(z)\to0$.
\end{proof}

The sign of $r_q$ on the real axis is the key to the next section: for real $x\in(0,1)$ the factor
$x(1-x)$ is positive, and if no twist $E_{\psi\chiK}$ vanishes at $x$, then $r_q(x)\ge0$.
